% Source printed pp. 1-4, through Figure 1.0.2.
\chapter{引言与动机}\label{src:1}
\begin{epigraph}
模空间理论不仅是一种绘制地图的方式, 还具有一个美妙的特点: 它有许多入口, 可以借助许多不同的技巧来发展. 当然, 有传统的代数几何, 但也有不变量理论, Teichmüller 理论等复解析方法, 整体拓扑方法, 以及在有限域上计数对象之类纯粹属于特征 $p$ 的方法. 这也是它魅力的一部分.
\epigraphauthor{David Mumford, \cite[前言]{Mum04}}
\end{epigraph}
模空间是作为数学中最基本的问题之一的解答而出现的: 
\begin{quote}
\textbf{分类问题: }能否对某一给定类型的数学对象的同构类进行分类?
\end{quote}
我们可能希望分类的对象有很多种: 
\begin{itemize}
\item 维数为 $k$ 的子空间 $V\subseteq\mathbb C^n$; 
\item 次数为 $d$ 的平面曲线 $C\subseteq\mathbb P^2$; 
\item 亏格为 $g$ 的曲线 $C$, 连同一个次数为 $d$ 的态射 $C\to\mathbb P^1$; 
\item 固定射影簇 $X$ 上的向量丛; 
\item 群的表示, 例如绝对 Galois 群 $\Gal(\overline{\mathbb Q}_p/\mathbb Q_p)$ 的表示, 亏格为 $g$ 的拓扑曲面 $\Sigma$ 的基本群 $\pi_1(\Sigma)$ 的表示, 或箭图的路径代数的表示.
\end{itemize}
我们主要关心以下两个例子, 全书将用它们来说明模的概念: 
\begin{enumerate}[(1)]
\item 亏格为 $g$ 的光滑(或更一般地, 稳定)射影曲线; 
\item 固定光滑曲线 $C$ 上的向量丛(更具体地说, 半稳定向量丛).
\end{enumerate}
\begin{quote}
\newterm{模空间}{moduli space}是一个空间, 其点与某一给定类型对象的同构类之间存在自然的双射.
\end{quote}
上面的关键词是``自然'', 你可能还不清楚这是什么意思. 事实上, 精确表述这一点, 正是发展模理论时的主要挑战之一. 毕竟, 任意两个正维的复流形或簇之间都存在双射, 因为它们都具有连续统的基数. 我们不希望模空间只是一些点的不交并, 每个同构类对应一个点. 我们需要为对象规定一种彼此接近的含义. 为使``自然双射''的概念形式化, 我们将引入对象族的概念, 并要求空间 $S$ 上的每个对象族都唯一对应于一个从 $S$ 到模空间的映射; 见第~\bookref{1.3}~节.

``空间''的结构取决于语境: 如果分类的是拓扑对象, 我们或许希望它具有拓扑空间结构; 如果分类的是解析对象, 我们或许希望它具有流形结构. 本书主要关注代数结构的分类, 希望模空间具有代数簇结构, 最好是射影簇结构.

\paragraph*{模的无处不在.} 一旦开始从模的角度看待空间, 一切似乎都像模空间!这在精确的意义下是正确的: 每个空间 $M$ 都是其点的模空间. 当然, 当存在具有几何意义的等价描述时, 情形就更有趣. 射影空间 $\mathbb P^1$ 可以看成 $\mathbb P^1$ 中所有点的集合(不太有趣), 也可以看成平面中过原点的直线的集合(更有趣), 或者 $\mathbb A^2$ 中支撑在原点, 长度为 $2$ 的子概形的集合(也很有趣), 又或者稳定椭圆曲线的同构类的集合(非常有趣). 引人入胜的是, 模空间常常能够在不同数学领域中获得等价的解释, 我们主要关心的两个例子正是如此: 
\begin{enumerate}[(1)]
\item $\mathcal M_g$ 是亏格为 $g$ 的光滑射影代数曲线的模, 也可以是亏格为 $g$ 的紧 Riemann 曲面的模, 或者固定的亏格为 $g$ 的紧定向拓扑曲面 $\Sigma_g$ 上按双全纯映射等同的复结构的模, 或者 $\Sigma_g$ 上按等距映射等同的双曲度量的模.
\item $\cBun_{r,d}^{\mathrm{ss}}(C)$ 是固定曲线 $C$ 上半稳定代数向量丛的模, 也可以是 $C$ 上带有平坦酉联络的全纯向量丛的模, 或者 $\pi_1(C)$ 的不可约酉表示的模.
\end{enumerate}
这使代数, 解析和拓扑方法之间产生了丰富的相互作用.

\paragraph*{离散的模与连续的模.} 根据对象的类型, 模空间可能是离散的, 也可能是连续的, 或者兼有二者.
\begin{itemize}
\item $\mathbb P^1$ 上线丛的模空间是离散集合 $\mathbb Z$: $\mathbb P^1$ 上每个线丛都同构于 $\mathcal O(n)$, 其中整数 $n\in\mathbb Z$ 唯一确定.
\item 参数化二次平面曲线 $C\subseteq\mathbb P^2$ 的模空间是连通空间 $\mathbb P^5$: 由 $a_0x^2+a_1xy+a_2xz+a_3y^2+a_4yz+a_5z^2$ 定义的曲线, 由点 $[a_0,\ldots,a_5]\in\mathbb P^5$ 唯一确定; 当平面曲线连续变化时(即改变系数 $a_i$), $\mathbb P^5$ 中相应的点也连续变化.
\item 对光滑曲线, 亏格 $g$ 是离散不变量. 固定亏格 $g$ 后, 光滑曲线的模空间 $M_g$ 是一个簇; 事实上, 它是一个不可约拟射影簇, 不过这样说已经远远超前了. 所有光滑曲线的模空间可以看成不交并 $\coprod_g M_g$.
\item 对光滑曲线 $C$ 上的向量丛, 秩 $r$ 和次数 $d$ 是离散不变量, 而秩为 $r$, 次数为 $d$ 的半稳定丛的模空间 $M_{r,d}^{\mathrm{ss}}(C)$ 是一个不可约射影簇.
\end{itemize}

\paragraph*{为什么研究模空间?} 模空间的性质能够告诉我们对象本身的性质. 对象的许多性质, 用模空间来表述最为合适. 例如, 为表达一般的亏格 $3$ 曲线能够用显式坐标系统参数化这一条件---一般的亏格 $3$ 曲线通过典范嵌入成为 $\mathbb P^2$ 中的一条四次平面曲线, 因而可以由空间 $\mathbb P(\Gamma(\mathbb P^2,\mathcal O(4)))\cong\mathbb P^{14}$ 中的点参数化---我们可以说模空间 $M_3$ 是单有理的, 即存在一个支配有理映射 $\mathbb P^{14}\dashrightarrow M_3$.

\paragraph*{模空间是否存在?} 在研究模空间的几何性质之前, 需要先知道它们存在. 这并非易事, 也是本书的主要目标. 我们建立模理论的基础, 以证明存在参数化亏格为 $g$ 的稳定曲线的射影模空间(定理~\hyperref[thm:A]{A}), 以及参数化固定光滑曲线上秩为 $r$, 次数为 $d$ 的半稳定向量丛的射影模空间(定理~\hyperref[thm:B]{B}).

代数对象的模空间常常作为代数簇存在, 这是一个令人惊叹的数学巧合. 它们的存在是模理论的起点. 模空间的美, 吸引力和优雅, 正是最初吸引我和无数其他人研究这一学科的原因.

\paragraph*{模的三分法.} 模理论中反复出现的一个主题, 是自同构群对模空间性质及其几何研究方法的影响. 根据自同构群的大小, 模理论有三种情形: 
\begin{enumerate}[(1)]
\item 没有自同构; 
\item 有限自同构; 
\item 无限自同构.
\end{enumerate}
在第三种情形中, 如果模叠的闭点---有时称为多重稳定对象---具有约化自同构群, 模空间的性质就尤其良好.
\begin{table}[htbp]
\centering\small
\begin{tabularx}{\textwidth}{@{}p{2.2cm}X X X@{}}
\toprule
自同构 & 无 & 有限 & 闭点处约化\\\midrule
空间类型 & 概形／代数空间 & Deligne--Mumford 叠 & 代数叠\\
定义性质 & Zariski／平展局部为仿射概形 & 平展局部为仿射概形 & 光滑局部为仿射概形\\
例子 & $\mathbb P^n$, $\mathrm{Gr}(q,n)$, $\mathrm{Hilb}$, $\mathrm{Quot}$ & $\mathcal M_g$ & $\cBun_{r,d}^{\mathrm{ss}}(C)$\\
商叠 $[X/G]$ & 作用自由 & 稳定子有限 & 闭轨道处稳定子约化\\
模空间的存在 & 细模空间 & 粗模空间 & 良模空间\\\bottomrule
\end{tabularx}
\caption{模的三分法.}\label{src:1.0.1}
\end{table}

\paragraph*{本章路线图.} 我们通过逐步赋予对象的模集合更丰富的结构, 说明本书采用这一方法的动机. 首先, 在第~\bookref{1.3}~节介绍对象族和函子观点, 然后在第~\bookref{1.4}~节发展群胚观点. 第~\bookref{1.5}~节说明引入平展拓扑的动机之后, 我们在第~\bookref{1.6}~节引入模叠, 将这些观点结合起来. 随后在第~\bookref{1.7}~节勾勒构造射影模空间的两种主要技巧.

\begin{figure}[htbp]
\centering
\begin{tikzcd}[math mode=false,ampersand replacement=\&,column sep=1.6em,row sep=2.2em,cells={nodes={align=center,font=\small}},every arrow/.append style={line width=.5pt}]
集合\arrow[d]\arrow[dr] \& \& 群胚\arrow[d]\\
拓扑空间\arrow[d] \& \shortstack{$\Sch$ 上的函子\\或预层}\arrow[d] \& $\Sch$ 上的预叠\arrow[d]\\
环化空间\arrow[d] \& $\Sch_{\text{ét}}$ 上的层\arrow[d] \& $\Sch_{\text{ét}}$ 上的叠\arrow[d]\arrow[dr]\\
概形\arrow[r,phantom,"$\subset$"] \& 代数空间\arrow[r,phantom,"$\subset$"] \& \shortstack{Deligne--Mumford\\叠}\arrow[r,phantom,"$\subset$"] \& 代数叠
\arrow[from=1-1,to=1-3,phantom,"$\subset$"]
\arrow[from=2-2,to=2-3,phantom,"$\subset$"]
\arrow[from=3-2,to=3-3,phantom,"$\subset$"]
\end{tikzcd}
\caption{集合与群胚逐步获得代数几何结构的示意图.}\label{src:1.0.2}
\end{figure}
